\chapter{Capital and property}
\label{ch:capital}

\section{Why the money has not come}

Cuba has been open to foreign investment since 1995 and has received very little
of it. The standard explanation is the embargo. The embargo is real and it is not
the main reason, and the evidence for that is simply that European, Canadian,
Chinese, Brazilian, Russian and Latin American capital faces no such prohibition
and has also stayed away.

The actual deterrents, in order of how much investors complain about them:

\emph{The employment agency.} A foreign venture does not hire its own workers. It
contracts them from a state agency, pays the agency in hard currency, and the
agency pays the worker in pesos, keeping the difference --- historically a
difference of more than an order of magnitude. The investor therefore pays a
first-world wage bill for a workforce receiving a fraction of it, has no control
over selection, cannot reward performance, and knows the arrangement is visible
to every employee. It is bad for the investor, worse for the worker, and its only
beneficiary is the state's foreign-currency account.

\emph{Approval by case.} Every project is negotiated and approved individually,
at ministerial and often at higher level, on timescales measured in years and
criteria that are not published. This is the definition of a discretionary
system, with the consequences Chapter~\ref{ch:integrity} describes.

\emph{No dispute resolution anyone trusts.} Contracts are justiciable in Cuban
courts, which have never ruled against the state on anything material. The
absence of credible arbitration is priced into every deal.

\emph{No repatriation guarantee that has held.} Dividends and capital may be
repatriated, in principle. In 2009 the state froze foreign firms' balances and in
several cases they remain unpaid. That is not a legal fact but a historical one,
and it is the one investors actually price.

\emph{Title III.} Since 2019, a company using property confiscated from an
American national is exposed to suit in United States courts. Since much of the
valuable commercial property in Cuba was confiscated from somebody, this is a
diffuse and unquantifiable liability attached to almost any large project, and
it will remain so until the claims are settled --- which is why
Design~\ref{d:debt}.2 puts them inside the debt restructuring.

The Mariel special development zone is the demonstration. It has excellent
infrastructure, a deep-water container port, a favourable tax regime, and after
more than a decade a small number of operating projects \citep{feinberg2016} and a total invested
capital far below what was projected. It fixed the tax rate and the customs
regime, which were not the problem, and kept the employment agency, the
case-by-case approval and the courts, which were.

\begin{design}{The investment regime}
\label{d:invest}
\begin{rules}
\rl{1}{Direct employment} The state employment agency monopoly is abolished.
Foreign-invested firms hire, pay, promote and dismiss their own staff under the
ordinary labour law, in the currency of Design~\ref{d:regime}. Applies to every
sector, including tourism (Design~\ref{d:visitor}.3) and the special zones.

\rl{2}{Approval by right, with a clock} Investment is permitted by right, subject
to a published negative list of restricted sectors and to ordinary regulatory
consents. Where an authorisation is genuinely required, a statutory decision
period applies and expiry of the period is a grant. Deemed consent is the single
most effective anti-corruption instrument in this book: an official who cannot
delay cannot extract.

\rl{3}{National treatment} Foreign and domestic investors, including Cuban
private firms and returning Cubans, face the same rules, taxes, access to inputs
and access to foreign exchange. The current arrangement, in which a foreign joint
venture may import and a Cuban private firm may not, is indefensible and is also
why Cuban businesses cannot grow.

\rl{4}{Arbitration} A standing offer of international arbitration under ICSID or
UNCITRAL rules, seated abroad, in the investment law itself and in every state
contract, plus a network of investment treaties. Cuba cannot borrow credibility;
it can rent it, and this is how.

\rl{5}{Repatriation} A guaranteed right to convert and repatriate profits and
capital, entrenched under Design~\ref{d:commit}, with the 2009 balances released
first as evidence.

\rl{6}{Publish the deals} Every investment authorisation, its terms, its
beneficial owners and any incentive granted, published. Per
Chapter~\ref{ch:integrity}, and it applies to military-linked entities without
exception.
\end{rules}
\end{design}

\section{Property at home, which comes first}

Here is a point that foreign-investment discussions consistently miss. Cuba's
largest property problem is not foreign; it is that Cubans own their homes and
cannot do anything with them.

The Urban Reform of 1960 transferred rental housing to its occupants and abolished
the residential rental market. It was the revolution's most popular and most
durable redistribution and it produced near-universal home ownership. It also
produced a housing stock in which nobody had a legal mechanism or a financial
incentive to invest, and which has been decaying for sixty-six years. Havana loses
buildings to collapse. The housing deficit runs to hundreds of thousands of
units.\unv{} Cubans hold what is, on paper, the country's single largest asset
class, and cannot borrow against it, cannot insure it properly, often cannot
prove title to it, and until 2011 could not sell it.

Chapter~\ref{ch:welfare} proposed a property tax as the natural local revenue.
That tax and this section are the same project: a cadastre, a registry, clean
title, and mortgageability turn a decaying liability into a financeable asset and
into the base of local government at the same time.

\begin{design}{Making Cuban property real}
\label{d:cadastre}
\begin{rules}
\rl{1}{Cadastre and registry} A complete national cadastre and public property
register, surveyed, digitised and open to inspection, with a defined programme
and a date. This is a large, dull, multi-year administrative project and it is a
precondition for the property tax, for mortgage lending, for
Design~\ref{d:tenure}, for building control, for insurance, and for the
settlement of claims. Nothing else in this chapter works without it.

\rl{2}{Title guaranteed} Registered title is guaranteed by the state and
defensible in the administrative court of Design~\ref{d:courts}.4, with a
compensation fund for registry error. A register nobody can rely on is a filing
cabinet.

\rl{3}{Regularisation} A time-limited programme to regularise informal
occupation, extensions and self-built housing on published criteria, favouring
the sitting occupier. A large share of Cuban housing is legally irregular after
six decades of improvisation, and a registry that excludes it excludes the
people it was meant to protect.

\rl{4}{Mortgage law} A mortgage statute with a workable, judicially supervised
enforcement procedure, plus consumer protections and a rehousing obligation on
foreclosure of a sole residence.

\rl{5}{Rental market} Residential letting legalised and regulated, with security
of tenure and published rent-review rules. Prohibiting rental did not make
housing affordable; it made it invisible, informal and unmaintained.
\end{rules}
\end{design}

\section{The claims}

Two bodies of claim, of very different size and character, and they have to be
settled together.

The \emph{certified United States claims} number 5,913, were certified at a
principal of \$1.9 billion, and at the statutory 6 per cent simple interest are
now valued at something over \$9 billion \citep{fcsc}. They are documented, adjudicated, and
held largely by corporations and their successors.

The \emph{uncertified claims} --- of Cuban Americans who were not United States
nationals at the time of taking, and of Cubans who never left and whose property
was taken under the same laws --- are far more numerous, far less documented,
emotionally far heavier, and potentially far larger in aggregate. They include
the houses left in the care of relatives in 1961 by people who expected to be
back in months.

Any settlement has to satisfy three constraints simultaneously. It must be
credible to the claimants, or Title III litigation continues and the country
stays uninvestable. It must be affordable to a bankrupt state. And it must not
displace a single Cuban from a home.

The central European experience is the guide, and it is a guide by negative
example. Restitution in kind --- returning the actual property --- was tried at
scale, most extensively in eastern Germany, where over two million claims were
lodged. It froze the property market for the better part of a decade, because no
purchaser and no lender would touch an asset that might be claimed; it consumed
enormous administrative and judicial capacity; and it delivered its benefits
overwhelmingly to people who had left, at the expense of people who had stayed.
Countries that compensated in money and paper instead --- Hungary's voucher
scheme, and the Baltic states in part --- got through it faster and cheaper.

\begin{design}{Settling the claims}
\label{d:claims}
\begin{rules}
\rl{1}{No restitution in kind} No occupied dwelling, no operating farm, no school
and no hospital is returned to a prior owner. Compensation only. This is stated
first, absolutely and without exception, because the fear of it is the single
biggest source of Cuban domestic resistance to normalisation, and removing that
fear is worth more than the assets involved.

\rl{2}{One process, one cap} A claims commission with a single filing window, a
published methodology, and an announced aggregate cap on total compensation,
inside the debt restructuring of Design~\ref{d:debt}. Claims are paid pro rata
within the cap. A settlement without a cap is not a settlement; it is an open
liability.

\rl{3}{Paid in paper} Compensation in the same long-dated instruments as the
restructured debt, with an option to convert into equity in privatisations under
Design~\ref{d:privat} or into investment in Cuba. Converting a claimant into an
investor is the best outcome available for both sides.

\rl{4}{Both bodies of claim} The uncertified Cuban and Cuban-American claims are
admitted to the same process on the same terms as the certified United States
claims. A settlement that pays foreign corporations and not Cuban families would
be politically impossible inside Cuba and morally indefensible anywhere.

\rl{5}{Finality} Participation extinguishes the claim, in all jurisdictions,
including under Title III. Finality is the entire product being purchased.

\rl{6}{Small claims fast} A simplified track with a low evidentiary threshold and
a flat payment for claims below a threshold, so that the large number of modest
family claims are resolved quickly rather than being ground through a process
designed for corporate assets.
\end{rules}
\end{design}

\section{Privatisation, and how not to create an oligarchy}

This is where transitions are lost, and Cuba has the advantage of thirty-five
years of other people's experience.

The lesson of the post-Soviet privatisations is not that privatisation is wrong.
It is that \emph{how} and \emph{in what order} determined everything. Mass voucher
schemes, in Russia and Czechoslovakia, dispersed nominal ownership among citizens
who promptly sold their vouchers for cash to the few actors with capital and
information, concentrating real ownership faster than a straightforward auction
would have. Loans-for-shares in Russia handed the largest assets to a handful of
people at derisory prices in exchange for financing a government. Poland, which
sold small enterprises early and fast, restructured medium ones slowly, and left
the largest for last under a functioning competition authority and a functioning
press, did better than anyone.

Cuba's specific risk is named in Chapter~\ref{ch:premises}: the incumbents
control the assets and the information, and a privatisation conducted before the
institutions exist hands them the economy with legal title attached. That outcome
is effectively irreversible, and it is the most likely failure mode of everything
in this book.

\begin{design}{Sequencing privatisation}
\label{d:privat}
\begin{rules}
\rl{1}{Small and many, first} Retail, catering, repair, transport, small
manufacturing and services --- thousands of units --- sold or leased early, in
open auctions, with a preference for the existing workers and for cooperatives.
These are the businesses the 1968 Revolutionary Offensive abolished, and
restoring them is both economically useful and politically popular.

\rl{2}{Large and few, last} Utilities, ports, telecommunications, mining and the
military conglomerate's core assets are not sold until the competition authority,
the sector regulators, the courts, the press and the beneficial-ownership
register are all functioning. Until then they are corporatised, audited and
published, but not transferred.

\rl{3}{Auctions, published} Every sale by open competitive auction with the
process, the bidders, the price and the beneficial owners published within thirty
days. No negotiated sales, no management buy-outs at book value, no strategic
investors selected on unpublished criteria.

\rl{4}{No vouchers} For the reasons above. Where wide ownership is the objective
it is pursued through worker and cooperative preference in clause~1 and through
the stabilisation fund holding minority stakes, not through paper distributed to
people with no means of valuing it.

\rl{5}{Beneficial ownership} No entity may acquire a privatised asset without
disclosing its ultimate beneficial owners to a public register. No anonymous
vehicles, no nominee structures, no exceptions for anybody.

\rl{6}{Cooling-off} Officials involved in a privatisation, and their households,
may not acquire an interest in the asset, work for the acquirer, or advise on the
transaction for a published period. Enforced by disclosure and by
Chapter~\ref{ch:integrity}'s asset declarations.

\rl{7}{Proceeds to the fund} Privatisation proceeds are rent under
Design~\ref{d:rentrule} and go to the stabilisation and investment fund. A
government that can spend privatisation receipts on the recurrent budget will
privatise for cash-flow reasons and at the wrong price.
\end{rules}
\end{design}

\section{GAESA}

The question nobody wants to write down, so it gets written down.

The armed forces' business conglomerate controls a large share of Cuba's
hard-currency economy --- hotels, retail chains, ports, logistics, remittance
processing, import and export --- and does not publish accounts. Its size is not
known outside itself, which is the first thing that has to change. It is
simultaneously the country's most competent commercial organisation, the largest
single misallocator of capital in the country per Chapter~\ref{ch:tourism}, and
the entity with the most to lose from every design in this book.

Three options, honestly assessed.

\emph{Expropriate and break up.} Politically impossible. The armed forces will
not permit it, they can prevent it, and a design that assumes otherwise is not a
design for Cuba. Attempting it is how a managed transition becomes a contested
one.

\emph{Leave it alone.} Guarantees the oligarchy outcome. A conglomerate with
privileged currency access, no published accounts and no competition, sitting
inside a liberalising economy, becomes its owner.

\emph{Normalise it.} Convert it into an ordinary state holding company, strip its
privileges, make it compete, publish its accounts, and let its constituent
businesses be sold off over a decade under Design~\ref{d:privat}. Slower and less
satisfying than the first, and it is the only one that is both achievable and
survivable.

\begin{design}{Normalising the military economy}
\label{d:gaesa}
\begin{rules}
\rl{1}{Civilian ownership} The conglomerate's holdings are transferred to a
civilian state holding company under the finance ministry. The armed forces
retain their budget, their pensions and their defence functions and cease to be a
commercial actor. Personnel keep their standing and their pensions under
Design~\ref{d:amnesty}.

\rl{2}{Published audited accounts} Full consolidated accounts, audited to
international standards by an external auditor, published annually, from the
first year. This is the clause that has to come first, because none of the others
can be assessed without it.

\rl{3}{No privileges} No preferential access to foreign currency, to import
licences, to land, to grid capacity, to state contracts or to regulatory
forbearance. Every clause of Design~\ref{d:invest} and Design~\ref{d:privat}
applies to it identically.

\rl{4}{Inside the budget} Its revenues and transfers appear in the single budget
of Design~\ref{d:fiscal}.1.

\rl{5}{Divested over a decade} Constituent businesses sold under
Design~\ref{d:privat}, small first, on a published schedule, with the proceeds to
the stabilisation fund and with the cooling-off rule of
Design~\ref{d:privat}.6 applied strictly to serving and former officers.

\rl{6}{The bargain is explicit} This clause exists to say plainly what
Chapter~\ref{ch:premises} argued: the surrender of commercial discretion is
exchanged for security, pensions, property and a legal route into ordinary
political and business life. That exchange should be negotiated openly and
written down, because a bargain left implicit is one that gets renegotiated at
every step, and Cuba cannot afford to renegotiate this one annually.
\end{rules}
\end{design}

\section{The diaspora as the first investor}

Finally, the observation that makes the rest of this chapter plausible.

Cuba is not going to attract large amounts of general international capital
quickly. It is a small, poor, defaulted, sanctioned market with no track record,
and the global competition for investment is intense. What it does have is two
million people abroad who know the country, speak the language, have relatives
there, and are already sending money in volumes that dwarf foreign direct
investment.

Remittances are consumption. The design objective is to convert a share of them
into investment, and the instruments are not exotic: dual nationality, the right
to own property and businesses on national treatment, a payment system that works
(Design~\ref{d:pay2}), a registry that gives a Miami investor confidence that the
title exists (Design~\ref{d:cadastre}), a court that a Cuban American can use
(Design~\ref{d:courts}.4), and a claims settlement that closes the past rather
than leaving it as an open wound (Design~\ref{d:claims}).

None of that requires the embargo to end. Cuban Americans may already travel,
send money and, within limits, support private Cuban businesses. The binding
constraints are on the Cuban side, and they are the same ones this chapter has
been listing.
